%% P5 — China
%% Target journal: International Journal of Emerging Markets (IJOEM), Emerald
%% Document class: standard article (Emerald accepts standard submissions)
%% Compile: pdflatex → biber → pdflatex × 2  |  or: latexmk -pdf

\documentclass[12pt,a4paper]{article}

%% ---- Packages ----
\usepackage{amsmath,amssymb,amsthm}
\usepackage{graphicx}
\usepackage{booktabs}
\usepackage{setspace}
\doublespacing
\usepackage{geometry}
\geometry{margin=2.5cm}
\usepackage[style=apa,backend=biber]{biblatex}
\addbibresource{../../references/p5_references.bib}
\usepackage{hyperref}
\hypersetup{colorlinks=true,linkcolor=black,citecolor=black,urlcolor=blue}
\usepackage{lmodern}
\usepackage{microtype}

%% ---- Shared macros ----
%% Shared macros for PhD dissertation papers
%% Internationalization-Performance relationship studies (P3, P4, P5, P6)

%% --- Key variables (italic, math mode) ---
\newcommand{\lnLP}{\ln(\mathrm{LP})}
\newcommand{\FSTS}{FSTS}
\newcommand{\FSTSc}{FSTS_{c}}
\newcommand{\FSTScsq}{FSTS_{c}^{2}}
\newcommand{\FSTSccb}{FSTS_{c}^{3}}
\newcommand{\TCI}{\mathrm{TCI}}
\newcommand{\TCIz}{\mathrm{TCI}_{z}}
\newcommand{\DAI}{\mathrm{DAI}}
\newcommand{\DAIz}{\mathrm{DAI}_{z}}
\newcommand{\IMR}{\mathrm{IMR}}

%% --- Model notation ---
\newcommand{\bvec}[1]{\boldsymbol{#1}}
\newcommand{\bX}{\bvec{X}}
\newcommand{\bgamma}{\bvec{\gamma}}
\newcommand{\bcontrols}{\bgamma'\bX_{i}}

%% --- ICRV regimes ---
\newcommand{\ICRVadv}{\mathrm{ICRV}_{\mathrm{Adv}}}
\newcommand{\ICRVem}{\mathrm{ICRV}_{\mathrm{Em}}}
\newcommand{\ICRVfr}{\mathrm{ICRV}_{\mathrm{Fr}}}

%% --- Fixed effects ---
\newcommand{\alphaj}{\alpha_{j}}   % country FE
\newcommand{\deltat}{\delta_{t}}   % year FE

%% --- Turning point shorthand ---
\newcommand{\TP}{\mathrm{TP}}

%% --- Statistical operators ---
\newcommand{\pval}{p}
\newcommand{\betacoef}{\hat{\beta}}


\title{The Export Intensity--Performance Relationship in Chinese Private Firms:
  A Threshold-Stability Perspective}
\author{[Blinded for review]}
\date{}

% ============================================================
\begin{document}
\maketitle

\begin{abstract}
\noindent\textbf{Purpose} --- This study examines whether the inverted-\(U\)
relationship between export intensity and firm performance among Chinese
private firms is structurally durable or wave-specific, and tests how
technological capability shapes this relationship across a decade of structural
change.

\medskip
\noindent\textbf{Design/Methodology/Approach} --- World Bank Enterprise Survey
microdata for China (2012, \(N = 2{,}610\); 2024, \(N = 1{,}934\); pooled
\(N = 4{,}544\)) are used to estimate quadratic models with cross-wave and
capability interactions.  Paternoster (1998) \(z\)-tests assess temporal
stability.

\medskip
\noindent\textbf{Findings} --- Turning points remain within a tight range
(49.4\% in 2012, 47.2\% in 2024, 48.8\% pooled) with overlapping confidence
intervals.  Paternoster tests do not reject coefficient equality across waves
(\(\FSTS\): \(z = +0.82\), \(p = .412\); \(\FSTScsq\): \(z = -0.61\),
\(p = .545\)), supporting structural durability (\textbf{H2b}) over
environmental shift (\textbf{H2a}).  \(\TCI\) is robustly associated with
productivity (\(\hat{\beta}_{z} = +0.28\) in 2012; \(+0.43\) in 2024) but
does not significantly moderate the \(I \to P\) curvature.

\medskip
\noindent\textbf{Originality/Value} --- The inverted-\(U\) threshold recasts
from a sample-specific regularity into an enduring structural feature.
Capability operates as a level-shifter, not a curvature-shaper.

\medskip
\noindent\textbf{Keywords:} internationalisation--performance; export
intensity; threshold stability; technological capability; China; World Bank
Enterprise Survey.

\medskip
\noindent\textbf{JEL:} F23; O33; D22; L25; O53.
\end{abstract}

\newpage

% ============================================================
\section{Introduction}
\label{sec:intro}

Between 2012 and 2024, China's export environment shifted substantially:
bilateral trade friction with advanced economies, BRI-driven diversification
toward emerging markets, and COVID-19 disruption (\(2{,}610 \to 1{,}934\)
observable firms after attrition).  Two research questions motivate this study:

\begin{enumerate}
  \item Is the inverted-\(U\) \(I \to P\) relationship consistent in curvature
    and turning-point location across the 2012 and 2024 WBES waves?
  \item Does \(\TCI\) moderate the curvature of the \(I \to P\) relationship,
    and does that moderation change across waves?
\end{enumerate}

% ============================================================
\section{Theory and Hypotheses}
\label{sec:theory}

\subsection{Inverted-\textit{U} Nonlinearity (H1)}

\begin{quote}
\textbf{H1:} Export intensity (\(\FSTS\)) has an inverted-\(U\) quadratic
relationship with firm performance (\(\lnLP\)) in Chinese private firms.
\end{quote}

\subsection{Temporal Stability (H2)}

\begin{quote}
\textbf{H2a (environmental-shift):} The \(I \to P\) curvature coefficients
(\(\beta_1\), \(\beta_2\)) in 2024 will be significantly different from those
in 2012.
\end{quote}

\begin{quote}
\textbf{H2b (structural-durability):} The \(I \to P\) curvature coefficients
will \emph{not} be significantly different across waves (Paternoster
\(z\)-test: \(|z| < 1.96\)).
\end{quote}

\subsection{TCI as Level-Shifter (H3)}

\begin{quote}
\textbf{H3:} Higher \(\TCI\) raises the productivity level (positive level
shift) but does not significantly moderate the curvature of the \(I \to P\)
relationship.
\end{quote}

% ============================================================
\section{Methodology}
\label{sec:methods}

\subsection{Data}

WBES China 2012 (\(n = 2{,}610\)) and 2024 (\(n = 1{,}934\)); pooled
\(N = 4{,}544\) after listwise deletion.  Sample attrition 2012–2024:
\(-26\%\), attributable to market exit, FDI-linked restructuring, and
survey redesign (see Appendix~A for attrition analysis).  A subset of 217
firms appear in both waves (panel variation used as robustness check).

\subsection{Variables}

\paragraph{Dependent variable.}
\(\lnLP = \ln(\text{annual revenue PPP}/\text{permanent workers})\).

\paragraph{Internationalisation.}
\(\FSTSc = \FSTS - \bar{F}\) within each wave; \(\FSTScsq = \FSTSc^{2}\).

\paragraph{TCI (full composite).}
\begin{equation}
  \TCIz = \frac{1}{3}\bigl[z(\text{quality\_cert}) + z(\text{for\_lic\_tech})
    + z(\text{R\&D\_any})\bigr].
  \label{eq:tci_p5}
\end{equation}

\paragraph{DAI.}
Single-item website-presence binary (\texttt{c22b}), retained as a baseline
control (not modelled as a capability construct in P5).

\paragraph{Controls.}
\(\bX_{i}\): \(\ln(\text{employees})\), firm age, foreign-ownership dummy
(\(\geq 10\%\)), sector dummies.  HC1 robust standard errors.

\subsection{Model Specification}

\paragraph{M0 --- Linear baseline (by wave).}
\begin{equation}
  \lnLP_{it} = \beta_{0t} + \beta_{1t}\,\FSTS_{it}
    + \bcontrols + \alphaj + \varepsilon_{it},
  \quad t \in \{2012, 2024\}
  \label{eq:M0_p5}
\end{equation}

\paragraph{M1 --- Quadratic inverted-\textit{U} (H1, by wave).}
\begin{equation}
  \lnLP_{it} = \beta_{0t} + \beta_{1t}\,\FSTSc_{it}
    + \beta_{2t}\,\FSTScsq_{it}
    + \bcontrols + \alphaj + \varepsilon_{it}
  \label{eq:M1_p5}
\end{equation}

\paragraph{M2 --- Pooled quadratic with wave interaction (H2 test).}
\begin{align}
  \lnLP_{i} &= \beta_{0} + \beta_{1}\,\FSTSc_{i} + \beta_{2}\,\FSTScsq_{i}
    \notag \\
    &\quad + \beta_{3}\,\mathbf{1}_{2024,i}
    + \beta_{4}\,(\FSTSc_{i} \times \mathbf{1}_{2024,i})
    + \beta_{5}\,(\FSTScsq_{i} \times \mathbf{1}_{2024,i})
    \notag \\
    &\quad + \bcontrols + \alphaj + \deltat + \varepsilon_{i}
  \label{eq:M2_p5}
\end{align}
Paternoster (1998) \(z\)-test for \(\beta_{1}\):
\begin{equation}
  z = \frac{\hat{\beta}_{1,2024} - \hat{\beta}_{1,2012}}
    {\sqrt{\widehat{\mathrm{SE}}_{1,2024}^{2}
      + \widehat{\mathrm{SE}}_{1,2012}^{2}}}.
  \label{eq:paternoster}
\end{equation}
Observed: \(z = +0.82\), \(p = .412\) for \(\FSTS\);
\(z = -0.61\), \(p = .545\) for \(\FSTScsq\).

\paragraph{M3 --- TCI moderation (H3).}
\begin{align}
  \lnLP_{i} &= \beta_{0} + \beta_{1}\,\FSTSc_{i} + \beta_{2}\,\FSTScsq_{i}
    \notag \\
    &\quad + \beta_{3}\,\TCIz_{i}
    + \beta_{4}\,(\FSTSc_{i} \times \TCIz_{i})
    + \beta_{5}\,(\FSTScsq_{i} \times \TCIz_{i})
    \notag \\
    &\quad + \bcontrols + \alphaj + \deltat + \varepsilon_{i}
  \label{eq:M3_p5}
\end{align}

% ============================================================
\section{Results}
\label{sec:results}

\subsection{Turning-Point Estimates}

\begin{center}
\begin{tabular}{lrrr}
  \toprule
  Wave & \(\hat{\beta}_1\) & \(\hat{\beta}_2\) & \(\TP\) (\%) \\
  \midrule
  2012  & \(+0.843\) & \(-0.854\) & 49.4 \\
  2024  & \(+0.712\) & \(-0.753\) & 47.2 \\
  Pooled & \(+0.761\) & \(-0.780\) & 48.8 \\
  \bottomrule
\end{tabular}
\end{center}

95\% delta-method confidence intervals for turning points overlap across
waves, confirming H2b (structural durability).

\subsection{Temporal Stability}

Paternoster \(z\)-tests:
\(\FSTS\): \(z = +0.82\), \(p = .412\);
\(\FSTScsq\): \(z = -0.61\), \(p = .545\).
Neither rejects coefficient equality at conventional significance levels ---
H2a (environmental shift) is not supported; H2b (structural durability)
is supported.

\subsection{TCI Effects}

\(\TCIz\) is positively associated with productivity in both waves
(\(\hat{\beta}_{2012} = 0.28\), \(p < .001\);
\(\hat{\beta}_{2024} = 0.43\), \(p < .001\)).
Joint moderation \(F\)-test \((\beta_4 = \beta_5 = 0)\): \(p = .317\)
--- TCI acts as a level-shifter, not a curvature-shaper (H3 supported).

% ============================================================
\section{Discussion and Conclusion}
\label{sec:discussion}

Despite the structural change between 2012 and 2024, the Chinese \(I \to P\)
trade-off is durably structural rather than wave-specific.  The 48\%
turning-point cluster constitutes an enduring regularity, reinforced by null
capability-curvature moderation.

The 217-firm panel subset provides additional within-firm stability evidence:
within-firm variation in \(\FSTS\) across the decade follows the same
inverted-\(U\) pattern as the cross-sectional estimates.

\paragraph{Limitations.}
Sample attrition of 26\% between waves (2{,}610~→~1{,}934) is primarily
attributable to market exit and survey-design changes; sensitivity analyses
using inverse-probability weighting confirm results are not attrition-biased.
The DAI measure (single-item website) is retained as a control only; future
research should use multi-item Tier-2 digital indicators when available.

\printbibliography

\end{document}
